\subsection{مصادر البيانات واختيار العينات وإعادة استخدامها}
\label{sec:datasets}
% Generated by scripts/build_dataset_overlap_evidence.py.
\newcommand{\DsExcerptPosts}{25}
\newcommand{\DsIndependentPosts}{10}
\newcommand{\DsDevelopmentReused}{7}
\newcommand{\DsDevelopmentAdditional}{3}
\newcommand{\DsArchivePanelPosts}{6}
\newcommand{\DsArchiveReused}{6}
\newcommand{\DsArchivePanelOverlap}{0}
\newcommand{\DsPromptPosts}{48}


يتبع اختيار البيانات السؤال الذي نريد اختباره. فمقارنة جودة التعليقات تحتاج إلى سياقات مشتركة وقاعدة واضحة لاختيار الأزواج المولّدة. ويضيف اختبار استخراج الرسالة سؤالًا عن مصدر نصها؛ إذ قد يختلف ضغط مقتطف من المصادر عن ضغط رسالة معيّنة بصورة مستقلة. أما اختبار إعادة بناء السياق، فيتطلب أن يحصل الطرفان على مدخلاتهما كلٌّ على حدة. تفسر هذه الفروق استخدام أكثر من مجموعة بيانات في التقييم.

يتيح الجدول~\ref{tab:dataset-map} مرجعًا مختصرًا لهذه المجموعات والاختبارات المساندة لها. تشمل صفوفه العشرة تجارب رسائل ومقارنات لتطوير تعليمات الكتابة واختبارات لتكرار التحضير وتشخيص النقل، ولا تمثل عشر مجموعات بيانات مستقلة. ويربط الشرح التالي بين قرارات الاختيار، بينما يحدد الجدول~\ref{tab:evidence-levels} النتائج المقابلة لها.

\begin{table}[!htbp]
\centering\small
\caption{مصادر السياق واختيار الحمولة في مجموعات التجارب والتحقق العشر. تختلف وحدة العد بين المنشورات وأزواج التعليقات والحالات التشخيصية.}
\label{tab:dataset-map}
\begin{RTL}
\begin{tabularx}{\columnwidth}{|A{0.23\columnwidth}|B|B|}
\hline
المجموعة & السياق وطريقة الاختيار & الحمولة أو علاقة إعادة الاستخدام \\ \hline
مقارنة الطريقتين في التجربة الأولى & مجموعة بحث من \NUM{167} منشورًا؛ تغطي أزواج جدول المقاييس \NUM{100} منشور. & تكرار التوليد لبعض المنشورات، وإتاحة مقتطف للمرجع بحسب ميزانية البتات. \\ \hline
تقييم جودة التعليقات & \NUM{1608} أزواج من \NUM{339} منشورًا، ثم اختيار زوج مؤهل لكل من \NUM{244} منشورًا. & تعيد مهام الحكم الخمس استخدام الأزواج نفسها؛ تخص أعداد نجاح التوليد التجربة الأوسع. \\ \hline
مقارنة تعليمات الكتابة & \NUM{\DsPromptPosts} منشورًا موزعًا بحسب نتائج عينة المحكّم. & تعيد الأنماط السبعة استخدام المنشورات ومواضع الرد، ولا تنقل الأنماط الضابطة دون موضوع رسائل. \\ \hline
اختبار العدد المحدود من التعليقات & أول \NUM{25} منشورًا في قائمة أرشيف مثبتة تضم \NUM{100} منشور. & حمولات معيّنة من \NUM{64} بتًا، وحد نقل من ثمانية تعليقات. \\ \hline
التحقق من إعادة بناء القوائم & تحضير المنشور نفسه بصورة مستقلة مرتين. & مقارنة النصوص والقوائم المرتبة؛ لا يمثل عينة رسائل منفصلة. \\ \hline
فحص سلسلة الرد & منشور من الاختبار المحدود، ثم تخطيط لأول \NUM{25} منشورًا منه. & حمولة من \NUM{64} بتًا وقائمة مرتبطة بالتعليق الأب؛ ليست عينة جديدة. \\ \hline
اختبارات النقل دون نموذج & مدخلات مضبوطة للتعليقات وتحديد الموضوع، مع حالات إتلاف متعمد. & فحص جمع البتات ورفض التتابعات المتضررة، وإعادة فحص حمولات مقتطفات المصادر. \\ \hline
استخراج مقتطفات المصادر & \NUM{\DsExcerptPosts} منشورًا إخباريًا بمدخلات محضّرة مثبتة. & رسائل مقتطفة من النصوص المشتركة لاختبار الضغط المعتمد على القاموس. \\ \hline
التطوير بحمولات مستقلة & \NUM{\DsIndependentPosts} منشورات: \NUM{\DsDevelopmentReused} معاد استخدامها و\NUM{\DsDevelopmentAdditional} ذات عنوان فقط. & رسائل مستقلة من ثمانية بايتات؛ جُرّبت سبع قبل التوقف. \\ \hline
إعادة البناء من الأرشيف & ستة منشورات لمقارنة الكتابة، وعينتا تطوير وتقييم منفصل من ستة منشورات لكل منهما. & يعيد التطوير استخدام الستة السابقة، وتستخدم العينة المنفصلة منشورات مختلفة، مع رسائل مستقلة. \\ \hline
\end{tabularx}
\end{RTL}
\end{table}

\subsubsection{مقارنة التعليقات في سياقات مشتركة}
تستند مقارنة الطريقتين في التجربة الأولى إلى مجموعة بحث من \NUM{167} منشورًا، تضم \NUM{139} منشورًا محضّرًا سابقًا و\NUM{28} منشورًا جرى البحث عنه حديثًا. حدّت حصص خدمة البحث من توسيع المجموعة، فأعادت التجربة استخدام منشورات بدل اعتبار كل تعليق مولّد سياقًا مصدريًا جديدًا. وأسهم كل من \NUM{35} منشورًا بعشر عينات. أما جدول المقاييس، فيضم في النهاية \NUM{304} أزواج مقبولة تغطي \NUM{100} منشور. تمثل هذه الأعداد مراحل متتابعة للاختيار، ولا تعبّر عن حجم واحد للعينة.

تأتي مجموعة تقييم جودة التعليقات من تجربة منفصلة، وتختار زوجًا مؤهلًا واحدًا لكل منشور. ويشترط القبول أن يكون تعليق \EN{ZLG}~\cite{lin2024zero} فريدًا داخل المجموعة المثبتة. يقلل الشرط تكرار نصوص المرجع في مهمة الحكم، لكنه يجعل العينة مختارة وفق شرط محدد، لا مسحوبة عشوائيًا. وتستخدم معايير الحكم الخمسة الأزواج نفسها، بينما تصف أعداد التوليد الأكبر التجربة التي أُخذت منها العينة.

تختلف طريقتا اختيار مجموعتي المقارنة. ولا تمثل إحداهما مجموعة تدريب والأخرى مجموعة اختبار، ولا نفترض خلوهما من المنشورات المشتركة. وتشترك التعليقات المتكررة للمنشور الواحد في السياق نفسه؛ لذلك تُجمع النتائج على مستوى المنشور حيث يُذكر ذلك.
ويختار فحص تطوير تعليمات الكتابة في القسم~\ref{sec:prompt-ablation} منشوراتٍ عددها \NUM{\DsPromptPosts} من عينة المحكّم وفق حالات الكشف المشتركة المعروفة. ويعيد استخدام سياقاتها لفحص مكونات الكتابة، ولا يزيد حجم عينة الجودة المستقلة.

\subsubsection{من جودة التعليق إلى استخراج الرسالة}
تفحص مقارنات الجودة النص الذي يقرؤه الناس. أما استخراج الرسالة، فيضيف شرطًا آخر: أن يستطيع المُستقبِل قراءة جميع الاختيارات التي تحمل رسالة واحدة. يفحص الاختبار التمهيدي وفحص سلسلة الرد السعةَ وتوافق الموضوع مع سياقه، بينما تختبر تجربتا المقتطفات والحمولات المستقلة رسائل كاملة تحت شروط ضغط مختلفة.

يستخدم اختبار العدد المحدود من التعليقات منشورات مؤرشفة من \EN{AskReddit} و\EN{NoStupidQuestions}، مع قائمة مثبتة من \NUM{100} منشور، تُستخدم أول \NUM{25} منها للاختبار التمهيدي. وتُثبت حمولات الـ\NUM{64} بتًا قبل التنفيذ. ويعيد فحص سلسلة الرد استخدام منشور من هذه المجموعة، ثم يفحص إمكان التخطيط على أول \NUM{25} منشورًا نفسها. لذلك يضيف تشخيصًا لسياقات موجودة، ولا يمثل عينة مستقلة أخرى.

تستخدم تجربة مقتطفات المصادر \NUM{\DsExcerptPosts} منشورًا إخباريًا محضّرًا ومثبتًا، مع نصوصه المصدرية وقائمة موضوعاته المرتبة. يقرأ البرنامج الملفات المحضّرة بترتيب ثابت، ويختار الرسائل من قاموس الضغط المشترك. يختبر ذلك نقل نص مرتبط بالمصدر بدقة، ويتيح للضاغط مادة مطابقة. وتغيّر عينة التطوير بحمولات مستقلة شرط اختيار الرسالة؛ إذ تضم \NUM{\DsIndependentPosts} منشورات، منها \NUM{\DsDevelopmentReused} ذات متن أُعيد استخدامها من تجربة المقتطفات المشخّصة، و\NUM{\DsDevelopmentAdditional} منشورات أرشيفية ذات عنوان فقط. وينتمي أحد المنشورات الإضافية أيضًا إلى اختبار العدد المحدود من التعليقات. أما الرسائل، فهي تتابعات قابلة للطباعة من ثمانية بايتات تُعيّن بحسب معرّف المنشور، بصورة مستقلة عن نص القاموس. لذا تشترك تجربتا التطوير والمقتطفات في بعض السياقات، ولا تُجمع نتائجهما بوصفها تكرارات مستقلة. دخلت المنشورات العشرة مرحلة التخطيط، ودخلت سبعة مرحلة التوليد قبل أن تترك قاعدة التوقف ثلاثة دون تجربة.

تجلب تجارب الأرشيف المنشورات والتعليقات من الأرشيف العام. تختار مقارنة الكتابة السابقة أول ستة منشورات مؤهلة نجح جلبها بحسب ترتيب الجلب المسجل. وتعيد مقارنة سياسات التطوير استخدام المنشورات الستة نفسها: اثنان من \EN{AskScience} وأربعة من \EN{technology}. أما عينة التقييم المنفصلة، فتضم ستة منشورات مختلفة، بواقع اثنين من كل من \EN{AskScience} و\EN{technology} و\EN{news}. يشترط القبول توافر \NUM{24} تعليقًا مؤرشفًا على الأقل، فتُستبعد المحادثات قليلة التعليقات. جُلبت أربعة منشورات من عينة التقييم حديثًا، بينما جُلب الاثنان الآخران سابقًا دون استخدامهما في التوليد. لذلك يشير عدم الاستخدام السابق إلى التوليد وتطوير السياسات، ولا يعني أن جميع المصادر نُزّلت حديثًا. وتُشتق الرسائل المستقلة القابلة للطباعة، المؤلفة من ثمانية بايتات، من معرّف المصدر وقيمة خاصة بكل تجربة. وتستخدم السياسات داخل كل عينة المنشورات والرسائل نفسها لإتاحة مقارنة مزدوجة؛ فلا تضيف كل سياسة عينة مصادر مستقلة.

\subsubsection{تحديد موضع الإخفاق في مسار الاتصال}
يقارن التحقق بتكرار التحضير استعلامات البحث والنصوص المجمعة وترتيب الموضوعات للمدخل نفسه. أما اختبارات النقل دون نموذج، فتستخدم مشاهدات مضبوطة للموضوعات، تشمل تتابعات متضررة، لعزل جمع البتات ومعالجة الأخطاء. وتمثل الرسائل الثلاث السليمة والحالات الأربع المتضررة في جدول النظرة العامة حالات تشخيصية، لا عينة من سبعة منشورات. وتعيد اختبارات نقل مقتطفات المصادر الخمسة والعشرون استخدام السياقات المحضّرة والرسائل المقابلة لها. وتضيف اختبارات المكونات الموضحة في القسم~\ref{sec:receiver-operation} مستوى آخر من التحقق، ولا تُضاف إلى أعداد عينات التجارب.

تساعد هذه الاختبارات مجتمعة على الفصل بين تعذر بناء قوائم متطابقة، والخطأ في تحديد موضوع التعليق، والإخفاق في جمع بتاته. وتُفسَّر كل نتيجة ضمن عينتها، بدل جمع أعداد المجموعات العشر في حجم عينة واحد.
